% Part: sets-functions-relations
% Chapter: infinite
% Section: dedekind
%
\documentclass[../../../include/open-logic-section]{subfiles}

\begin{document}

\olfileid[tr]{sfr}{infinite}{dedekind}
\olsection{Dedekind Cebirleri}

Doğal sayıların yalnızca sonsuz olmasını değil, belirli (cebirsel)
özelliklere de sahip olmasını isteriz: toplama, çarpma ve benzeri işlemler
altında iyi davranmaları gerekir.

Dedekind'in düşüncesi, \emph{ardıl fonksiyonu} fikrini temel almak ve
ardından sayıları aşağıdaki özelliklere sahip nesneler olarak
nitelemekti:
\begin{enumerate}
	\item Hiçbir sayının ardılı olmayan bir $0$ sayısı vardır
	\\yani $0 \notin \ran{s}$
	\\yani $\forall x\ s(x) \neq 0$
	\item Farklı sayıların ardılları da farklıdır
	\\yani $s$ bir !!a{injection}dur
	\\yani $\forall x \forall y (s(x) = s(y) \lif x = y)$
	\item\ollabel{repeatedapplication} Her sayı, ardıl fonksiyonunun $0$'a
	tekrar tekrar uygulanmasıyla elde edilir.
\end{enumerate}
İlk iki koşulu birinci derece mantıkla ele almak kolaydır (yukarıya
bakın). Fakat yalnızca birinci derece mantığı kullanarak
\olref{repeatedapplication}'ı ele alamayız. Dedekind'in çığır açan adımı,
\olref{repeatedapplication} koşulunu küme teorisi terimleriyle şöyle yeniden
ifade etmekti:
\begin{enumerate}
	\item[3$'$.] Doğal sayılar, \emph{ardıl fonksiyonu altında kapalı} olan
	en küçük kümedir: yani $s$'yi kümenin herhangi bir !!{element}ına
	uygularsak yine kümenin bir !!{element}ını elde ederiz.
\end{enumerate}
Ancak bunu ağır ağır açmamız gerekecek.

\begin{defn}\ollabel{Closure}
	Her $A$ kümesi, her $f\colon A\to A$ fonksiyonu ve her $X\subseteq A$ için, $X$'e ancak
	ve ancak $(\forall x\in X)f(x)\in X$ ise \emph{$f$ altında kapalı} denir.
	Şimdi her $o\in A$ için şöyle tanımlayalım:
	$$\closureofunder{f}{o} = \bigcap\Setabs{X}{X\subseteq A,\ o \in X
\text{ ve $X$, $f$ altında kapalıdır}}$$
\end{defn}

Böylece $\closureofunder{f}{o}$, $o$'yu !!a{element} olarak içeren ve $f$
altında kapalı olan bütün $A$ alt kümelerinin kesişimidir. Dolayısıyla sezgisel
olarak $\closureofunder{f}{o}$, $o$'yu !!a{element} olarak içeren \emph{en
küçük} $f$ altında kapalı kümedir. Sıradaki sonuç bu sezgisel düşünceyi
kesinleştirir:
\begin{lem}\ollabel{closureproperties}
	Her $A$ kümesi, her $f\colon A\to A$ fonksiyonu ve her $o\in A$ için:
	\begin{enumerate}
		\item\ollabel{closurehaselem} $o \in \closureofunder{f}{o}$; ve
		\item\ollabel{closureclosed} $\closureofunder{f}{o}$, $f$ altında
		kapalıdır; ve
		\item\ollabel{closuresmallest} $X\subseteq A$, $f$ altında kapalı ve
		$o\in X$ ise $\closureofunder{f}{o} \subseteq X$.
	\end{enumerate}
\end{lem}

\begin{proof}
$o$'yu !!a{element} olarak içeren en az bir $f$ altında kapalı küme vardır;
bu küme $\ran{f}\cup\{o\}$'dur. Dolayısıyla bu tür kümelerin \emph{tümünün}
kesişimi olan $\closureofunder{f}{o}$ vardır. Şimdi
\olref{closurehaselem}--\olref{closuresmallest}'ı denetlemeliyiz.

\olref{closurehaselem} için: $\closureofunder{f}{o}$, hepsi $o$'yu
!!a{element} olarak içeren kümelerin kesişimi olduğundan
$o\in\closureofunder{f}{o}$.

\olref{closureclosed} için: $x\in\closureofunder{f}{o}$ olduğunu varsayın.
Öyleyse $o\in X$ ve $X$, $f$ altında kapalı olduğunda $x\in X$ olur; $X$,
$f$ altında kapalı olduğu için de $f(x)\in X$ olur. Dolayısıyla
$f(x)\in\closureofunder{f}{o}$.

\olref{closuresmallest} için: genel olarak $X\in C$ ise
$\bigcap C\subseteq X$ olur.
\end{proof}

Bunu kullanarak şunu söyleyebiliriz:

\begin{defn}
Bir \emph{Dedekind cebiri}, bir $A$ kümesi, bir $f\colon A\to A$ fonksiyonu
ve aşağıdakileri sağlayan bir $o\in A$'dan oluşur:
	\begin{enumerate}
		\item \ollabel{ded:proper} $o \notin \ran{f}$
		\item \ollabel{ded:injection} $f$ bir !!a{injection}dur
		\item \ollabel{ded:closure} $A = \closureofunder{f}{o}$
	\end{enumerate}
\end{defn}

$A=\closureofunder{f}{o}$ olduğundan önceki sonucumuz, $A$'nın $o$'yu
!!a{element} olarak içeren en küçük $f$ altında kapalı küme olduğunu söyler.
Bir Dedekind cebirinin Dedekind-sonsuz olduğu açıktır; tanımın
\olref{ded:proper} ve \olref{ded:injection} maddelerine bakmak yeterlidir.
Ancak daha heyecan verici olgu, her Dedekind-sonsuz kümenin bir Dedekind
cebirine dönüştürülebilmesidir.

\begin{thm}\ollabel{thm:DedekindInfiniteAlgebra}
Dedekind-sonsuz bir küme varsa bir Dedekind cebiri vardır.
\end{thm}

\begin{proof}
$D$ Dedekind-sonsuz olsun. Öyleyse $g\colon D\to D$ bire bir fonksiyonu ve
bir $o\in D\setminus\ran{g}$ vardır. Şimdi
$A=\closureofunder{g}{o}$ olsun; \olref{closureproperties}'a göre $A$
vardır ve $o\in A$'dır. $f=\funrestrictionto{g}{A}$ olsun. $A,f,o$'nun bir
Dedekind cebiri oluşturduğunu göstereceğiz.

\olref{ded:proper} için: $o\notin\ran{g}$ ve $\ran{f}\subseteq\ran{g}$
olduğundan $o\notin\ran{f}$.

\olref{ded:injection} için: $g$, $D$ üzerinde bire birdir; öyleyse
$f\subseteq g$ de bire bir olmalıdır.

\olref{ded:closure} için: \olref{closureproperties}'a göre $A$, $g$ altında
kapalıdır; dolayısıyla özellikle $f$ altında da kapalıdır. Bu yüzden yine
\olref{closureproperties}'a göre $\closureofunder{f}{o}\subseteq A$.
Ayrıca $\closureofunder{f}{o}$, $f$ altında kapalı ve
$f=\funrestrictionto{g}{A}$ olduğundan $g$ altında da kapalıdır. Dolayısıyla
\olref{closureproperties}'a göre
$A\subseteq\closureofunder{f}{o}$.
\end{proof}

\end{document}
